\label{sec:limits}
\begin{itemize}
\item \textbf{Scale.} One model (\rhval{const/n_params} parameters), one dataset, one budget of \rhval{const/steps} steps (far from convergence). Rankings between schedules are known to change with training length, and a short budget favours schedules that reach low loss quickly.
\item \textbf{Seeds and statistics.} Only \rhval{const/n_seeds} seeds per cell, so Welch/paired $p$-values rest on $n=3$ and the standard deviations are themselves noisy; many comparisons are decided by effect sizes that are large relative to the seed spread, but small differences (such as between cosine at neighbouring peak rates) are not resolved. We ran many comparisons without multiplicity correction.
\item \textbf{Untuned secondary choices.} Warmup length, cosine floor, step positions and factor, Adam betas, weight decay and the Schedule-Free configuration were fixed a priori, not tuned. A different Schedule-Free warmup or weighting could change its sensitivity; we cannot say whether its poor low-LR behaviour is intrinsic or due to our choices (the interplay of the squared-LR weights and the 100-step warmup may matter).
\item \textbf{Reimplementation.} All four systems, including Schedule-Free AdamW, are written in \texttt{method/run.py} from the published description; they were not checked against the authors' reference implementation numerically.
\item \textbf{Confounded baseline specification.} As the seed asked, constant and step decay have no warmup in the main table. We added a warmup ablation to separate the effects, but the fair-warmup runs cover only some cells (both only at the three main rates), and sensitivity of those variants is over three rates.
\item \textbf{Sensitivity metric.} A range over a coarse five-point log grid whose position relative to each optimum matters; Schedule-Free's best rate is at or above the upper end of the main grid. We did not measure sensitivity to other hyperparameters or interactions of peak LR with the decay horizon.
\item \textbf{Single metric.} We report validation loss on 256 fixed windows after the last step only; we did not evaluate at equal wall-clock (Schedule-Free's extra state costs memory and some time), generalisation gaps, or sample quality, and no test set was held out because no model selection was done on it. Learning-rate curves in Fig.~\ref{fig:curves} are single-seed.
\item \textbf{Mechanism.} We isolated warmup as a factor; we did not test why it helps.
\end{itemize}
